\section{Explainability}
\label{sec:xai}

A high-PR-AUC model is operationally useful only if its predictions can be interrogated. We attribute every test-set prediction to the features that produced it, and then aggregate those attributions in three views: a global view (which features matter on average), a local view (why this single stop is flagged), and an ablation view (what the model loses when a whole feature group is removed).

\subsection{Setup}

We pick the strongest tabular model per scenario by test PR-AUC -- XGBoost wins both -- and run TreeSHAP~\cite{lundberg2017unified,lundberg2020local} on a stratified \num{270644}-row sample (10\% of the \num{2706443}-row test set, fixed seed~42). The resulting SHAP value tensor of shape $(\num{270644} \times F)$ is persisted to disk as \texttt{shap\_values.npy} immediately after compute, so any subsequent figure derives from the same numbers. The lag-feature ablation refits and predicts on the full training and test sets (not the SHAP subsample). The full Phase~4 pipeline -- both scenarios, the lag-feature ablation retrain, and 14 figures -- runs in roughly 24 minutes.

\subsection{Global feature importance}

The global view averages the absolute SHAP value for each feature across the full test set. \Cref{fig:global-bar} shows the top 20 features for each scenario.

\begin{figure}[H]
    \centering
    \begin{subfigure}{0.48\textwidth}
        \centering
        \includegraphics[width=\linewidth]{stop_level/xai/global_summary_bar_A.png}
        \caption{Scenario A (pre-departure).}
    \end{subfigure}\hfill
    \begin{subfigure}{0.48\textwidth}
        \centering
        \includegraphics[width=\linewidth]{stop_level/xai/global_summary_bar_B.png}
        \caption{Scenario B (inflight).}
    \end{subfigure}
    \caption{Top-20 global feature importance, measured as mean $|\text{SHAP}|$ per feature. In the pre-departure scenario the dominant signal is \texttt{train\_station\_lag7\_rate}: the seven-day rolling reliability of \emph{this} train at \emph{this} station. In the inflight scenario \texttt{prev\_stop\_actual\_delay} dominates by a factor of roughly 3.7 over the next-strongest feature.}
    \label{fig:global-bar}
\end{figure}

The pre-departure model leans most on the rolling reliability of a specific (train, station) pair: \texttt{train\_station\_lag7\_rate} reaches a mean $|\text{SHAP}|$ of \num{1.061}, ahead of \texttt{avg\_historical\_delay} (\num{0.396}), \texttt{stop\_order} (\num{0.273}), \texttt{station\_lag7\_rate} (\num{0.194}) and \texttt{temperature} (\num{0.165}). The inflight model is dominated by a single feature: \texttt{prev\_stop\_actual\_delay} reaches mean $|\text{SHAP}|$ of \num{2.464}, nearly 3.7 times the second-strongest feature (\texttt{stop\_order}, \num{0.674}).

\Cref{fig:beeswarm} shows per-row SHAP distributions. Each dot is one test stop; horizontal position is the SHAP value; colour is the feature value. The scenario-B \texttt{prev\_stop\_actual\_delay} row spans roughly $-2$ to $+9$ SHAP units, and the high (red) values consistently push the prediction toward ``disrupted'' -- the textbook ``inflight cliff''.

\begin{figure}[H]
    \centering
    \begin{subfigure}{0.48\textwidth}
        \centering
        \includegraphics[width=\linewidth]{stop_level/xai/global_summary_beeswarm_A.png}
        \caption{Scenario A.}
    \end{subfigure}\hfill
    \begin{subfigure}{0.48\textwidth}
        \centering
        \includegraphics[width=\linewidth]{stop_level/xai/global_summary_beeswarm_B.png}
        \caption{Scenario B.}
    \end{subfigure}
    \caption{Per-row SHAP distributions (beeswarm). Each dot is one test stop. Horizontal position is the SHAP value, i.e.\ the feature's contribution to the log-odds of disruption. Colour encodes feature value (red high, blue low). The dense red cluster on the right of the scenario-B \texttt{prev\_stop\_actual\_delay} row is the inflight cliff.}
    \label{fig:beeswarm}
\end{figure}

\subsection{Lag-feature ablation}

To quantify how much of the test-set signal lives in the lag and inflight feature groups, we re-fit the winning model after removing every \texttt{*\_lag*} feature \emph{and} the three inflight features. The remaining feature set is purely static, weather and position. \Cref{tab:lag-ablation} shows the result.

\begin{table}[H]
    \centering
    \caption{Lag-feature ablation. ``PR-AUC (full)'' is the test-set score with the full feature set; ``PR-AUC (no lag)'' removes every lag feature and every inflight feature simultaneously. The two scenarios collapse to the same PR-AUC of \num{0.350} once those groups are gone, so the entire A$\to$B uplift is carried by lag and inflight features.}
    \label{tab:lag-ablation}
    \small
    \begin{tabular}{l S[table-format=1.3] S[table-format=1.3] S[table-format=1.3]}
        \toprule
        Scenario & {PR-AUC (full)} & {PR-AUC (no lag)} & {$\Delta$} \\
        \midrule
        A & 0.533 & 0.350 & 0.183 \\
        B & 0.940 & 0.350 & 0.590 \\
        \bottomrule
    \end{tabular}
\end{table}

\Cref{fig:lag-ablation} visualises the drop. With lag and inflight features removed, both scenarios converge to PR-AUC \num{0.350}, which is barely better than the disruption rate alone. Lag and inflight features carry essentially the whole A$\to$B uplift.

\begin{figure}[H]
    \centering
    \includegraphics[width=0.5\textwidth]{stop_level/xai/lag_ablation.png}
    \caption{Lag-feature ablation, scenario~B. The full model reaches PR-AUC \num{0.940}; with lag and inflight features removed the same model reaches only \num{0.350}, a drop of \num{0.590}.}
    \label{fig:lag-ablation}
\end{figure}

\subsection{Local explanations}

To complement the global view we pick four prototypical test rows: the highest-confidence true positive, the worst false negative, a confident false positive, and the row whose prediction flipped most dramatically between scenarios A and B. We then plot a SHAP waterfall for each (\cref{fig:waterfalls}; the false-positive and false-negative versions are in \cref{fig:appendix-waterfalls}). The true-positive plot, for service \texttt{FI\_24\_2024-06-06} at station \texttt{FI\_HKH}, shows \texttt{prev\_stop\_actual\_delay} contributing $+8.34$ SHAP units on its own -- more than every other feature combined. The model predicts disruption with $p=1.000$.

\begin{figure}[H]
    \centering
    \begin{subfigure}{0.45\textwidth}
        \centering
        \includegraphics[width=\linewidth]{stop_level/xai/local_waterfall_tp.png}
        \caption{Highest-confidence true positive.}
    \end{subfigure}\hfill
    \begin{subfigure}{0.45\textwidth}
        \centering
        \includegraphics[width=\linewidth]{stop_level/xai/local_waterfall_AB_flip.png}
        \caption{A$\leftrightarrow$B prediction flip.}
    \end{subfigure}
    \caption{SHAP waterfalls for two prototypical test rows. Bars push the prediction left (toward ``not disrupted'') or right (toward ``disrupted''). The TP plot is dominated by a single feature; the flip plot shows several features changing sign between scenarios.}
    \label{fig:waterfalls}
\end{figure}

\subsection{Scenario uplift across stop positions}

\Cref{fig:uplift} plots the per-row scenario-uplift $P_B - P_A$ against the normalised stop position. The mean uplift is $-\num{0.128}$: scenario~B is generally \emph{more} selective, not less. Scenario~A over-predicts; scenario~B uses the inflight signal to rule out many of those over-predictions. The uplift is more negative late in the route (\texttt{position\_norm} above \num{0.75}), where the inflight evidence has the most history to work with.

\begin{figure}[H]
    \centering
    \includegraphics[width=0.92\textwidth]{stop_level/xai/scenario_uplift_map.png}
    \caption{Per-row scenario uplift $P_B - P_A$ against \texttt{position\_norm}. Negative uplift means scenario~B is more selective than scenario~A. The mean uplift is $-\num{0.128}$ and the effect strengthens late in the route.}
    \label{fig:uplift}
\end{figure}
